\chapter{La topología de un subespacio}\label{chap:04}
\section{El espacio ambiente forma parte de la afirmación}
Decir que un conjunto es abierto o cerrado requiere indicar dentro de qué espacio se trabaja. El intervalo $[0,1]$ es cerrado en $\R$, pero es abierto dentro de sí mismo. Esta no es una contradicción: las bolas disponibles cambian al restringir el universo. La lectura 4 de \cite{valencia} estudia precisamente esta distinción. Usaremos las propiedades de bolas, abiertos y clausuras de los capítulos~\ref{chap:02} y \ref{chap:03}.

\begin{definicion}\index{subespacio!metrico@métrico}\label{ch04:subespacio}
Si $(X,d)$ es métrico y $Y\subset X$, la restricción $d|_{Y\times Y}$ es la métrica inducida en $Y$. El espacio resultante se llama subespacio métrico.
\end{definicion}
La restricción es no negativa, simétrica y separa puntos porque conserva los valores de $d$. Para $u,v,w\in Y$, la desigualdad triangular ya es válida en $X$ y, por tanto, en $Y$. Esto prueba todos los axiomas, sin ninguna hipótesis de que $Y$ sea abierto o cerrado.

Indicaremos el ambiente mediante subíndices cuando pueda haber confusión. Para $y\in Y$,
\begin{equation}\label{ch04:bolas}
B_Y(y,r)=Y\cap B_X(y,r).
\end{equation}
La identidad se deduce porque ambos miembros exigen simultáneamente pertenecer a $Y$ y tener distancia a $y$ menor que $r$. También vale con bolas cerradas y esferas.

\begin{proposicion}\label{ch04:vecindades}
Un conjunto $V\subset Y$ es vecindad de $y\in Y$ en $Y$ si y solo si existe una vecindad $W$ de $y$ en $X$ tal que $V=W\cap Y$.
\end{proposicion}
\begin{proof}
Si $B_Y(y,r)\subset V$, definamos $W=V\cup B_X(y,r)$. Es vecindad de $y$ en $X$ porque contiene la bola ambiental. Además,
\[
W\cap Y=(V\cap Y)\cup(B_X(y,r)\cap Y)
=V\cup B_Y(y,r)=V.
\]
La última igualdad usa $B_Y(y,r)\subset V$. Es una unión, no una intersección, en este cálculo distributivo. En sentido inverso, si $V=W\cap Y$ y $B_X(y,r)\subset W$, entonces $B_Y(y,r)=B_X(y,r)\cap Y\subset V$.
\end{proof}

\begin{teorema}\label{ch04:relativa}\index{topologia@topología!relativa}
Para $A\subset Y\subset X$:
\begin{enumerate}[label=(\roman*)]
\item $A$ es abierto en $Y$ si y solo si $A=Y\cap G$ para algún abierto $G$ de $X$.
\item $A$ es cerrado en $Y$ si y solo si $A=Y\cap F$ para algún cerrado $F$ de $X$.
\end{enumerate}
\end{teorema}
\begin{proof}
\begin{enumerate}[label=(\roman*)]
\item Si $A$ es abierto en $Y$, para cada $a\in A$ elegimos $r_a>0$ con $B_Y(a,r_a)\subset A$. Pongamos $G=\bigcup_{a\in A}B_X(a,r_a)$, abierto en $X$. Entonces
\[
Y\cap G=\bigcup_{a\in A}B_Y(a,r_a)=A,
\]
pues cada bola está en $A$ y cada $a$ pertenece a su bola. Si $A$ es vacío, se toma $G=\varnothing$. Recíprocamente, si $A=Y\cap G$ y $a\in A$, la apertura de $G$ proporciona una bola $B_X(a,r)\subset G$; su intersección con $Y$ está en $A$ y demuestra apertura relativa.
\item $A$ es cerrado en $Y$ si y solo si $Y\setminus A$ es abierto en $Y$. Por (i), esto equivale a $Y\setminus A=Y\cap G$ para algún abierto $G$ en $X$. Tomando complementos relativos, $A=Y\cap(X\setminus G)$, donde $X\setminus G$ es cerrado en $X$. En sentido inverso, si $A=Y\cap F$ con $F$ cerrado, $Y\setminus A=Y\cap(X\setminus F)$ es abierto en $Y$ por (i).
\end{enumerate}
\end{proof}

\begin{proposicion}\label{ch04:clausura}
Si $A\subset Y$, entonces $\cl A^{\,Y}=Y\cap\cl A^{\,X}$, donde los superíndices indican el espacio de clausura. En consecuencia, $A$ es cerrado en $Y$ si y solo si $A=Y\cap\cl A^{\,X}$.
\end{proposicion}
\begin{proof}
Para $y\in Y$, la condición de pertenecer a la clausura relativa es que $B_Y(y,r)\cap A\ne\varnothing$ para todo $r>0$. Como $A\subset Y$, esta intersección es exactamente $B_X(y,r)\cap A$. La condición equivale, por tanto, a $y\in\cl A^{\,X}$. Para ser cerrado, $A$ debe coincidir con su clausura relativa.
\end{proof}
La clausura ambiental puede añadir puntos que no están disponibles en el subespacio. En $Y=(0,1)$, el conjunto $A=Y$ es cerrado relativamente, pero su clausura en $\R$ es $[0,1]$.

\begin{observacion}[Convergencia relativa]\label{ch04:convergencia}
Una sucesión de puntos de $Y$ converge a $y\in Y$ en el subespacio si y solo si converge a $y$ en $X$: las distancias numéricas son idénticas. Sin embargo, puede converger en $X$ a un punto exterior a $Y$, en cuyo caso no tiene ese límite en el subespacio. Por ejemplo, $1/n$ converge a $0$ en $\R$, pero no converge en $(0,1]$.
\end{observacion}

\section{Ejercicios resueltos}
\begin{ejercicio}[Lectura 4, ejercicio 1]
Demostrar que los abiertos y cerrados relativos son intersecciones con abiertos y cerrados del espacio ambiente.
\end{ejercicio}
\begin{solucion}
La demostración completa es el teorema~\ref{ch04:relativa}. Su construcción explícita, útil al aplicar el resultado, consiste en elegir para cada $a$ del abierto relativo una bola $B_Y(a,r_a)$ contenida en él y tomar en $X$ la unión $G=\bigcup_aB_X(a,r_a)$. Así se obtiene el abierto ambiental requerido. Para un cerrado relativo $A$, se aplica esa construcción a $Y\setminus A$ y se toma el complemento ambiental del abierto obtenido. Ambas direcciones y el caso vacío quedaron comprobados allí.
\end{solucion}

\begin{ejercicio}[Lectura 4, ejercicio 2]
Sea $A\subset Y\subset X$. Probar: (i) si $A$ es abierto en $Y$ y $Y$ abierto en $X$, entonces $A$ es abierto en $X$; (ii) si $A$ es cerrado en $Y$ y $Y$ cerrado en $X$, entonces $A$ es cerrado en $X$.
\end{ejercicio}
\begin{solucion}
\begin{enumerate}[label=(\roman*)]
\item Por el teorema~\ref{ch04:relativa}, $A=G\cap Y$ para algún abierto $G$ de $X$. Ambos factores son abiertos en $X$ y una intersección finita de abiertos es abierta.
\item Ahora $A=F\cap Y$ para algún cerrado $F$ de $X$. La intersección de cerrados es cerrada en $X$.
\end{enumerate}
Las hipótesis sobre $Y$ no son superfluas: $[0,1]$ es abierto en sí mismo y no en $\R$; $(0,1)$ es cerrado en sí mismo y no en $\R$.
\end{solucion}

\begin{ejercicio}[Lectura 4, ejercicio 3]
Considerar $\Q\subset\R$ con la métrica usual. Para un intervalo abierto relativo no vacío y acotado $I=(a,b)\cap\Q$, con $a<b$, demostrar que $I$ es cerrado en $\Q$ si y solo si ambos extremos de su clausura en $\R$ son irracionales. Precisar el caso de intervalos no acotados.
\end{ejercicio}
\begin{solucion}
La densidad de $\Q$ en $\R$ implica $\cl I^{\,\R}=[a,b]$. En detalle, todo punto de $[a,b]$ puede aproximarse por racionales de $(a,b)$; un punto fuera de $[a,b]$ tiene una bola que no intersecta al intervalo. Por la proposición~\ref{ch04:clausura},
\[
\cl I^{\,\Q}=[a,b]\cap\Q.
\]
Este conjunto coincide con $(a,b)\cap\Q$ exactamente cuando ni $a$ ni $b$ son racionales. Si un extremo es racional, pertenece a la clausura relativa y falta en $I$, por lo que este no es cerrado.

Para $(-\infty,b)\cap\Q$, solo importa el extremo finito $b$: el conjunto es cerrado en $\Q$ exactamente cuando $b\notin\Q$. Para $(a,\infty)\cap\Q$ ocurre lo mismo con $a$. El intervalo $\Q$ entero es abierto y cerrado. Así, en la versión sin acotación debe decirse «todos los extremos finitos son irracionales»; $\pm\infty$ no son números irracionales. El intervalo vacío también es abierto y cerrado y se trata por separado.
\end{solucion}

\begin{ejercicio}[Lectura 4, ejercicio 4]
Para $A\subset Y\subset X$, demostrar que el derivado relativo $A'_Y$ está contenido en $A'_X$, y dar un ejemplo de inclusión estricta.
\end{ejercicio}
\begin{solucion}
En realidad se tiene la identidad $A'_Y=Y\cap A'_X$. Para $y\in Y$, la condición $B_Y(y,r)\cap(A\setminus\{y\})\ne\varnothing$ equivale a $B_X(y,r)\cap(A\setminus\{y\})\ne\varnothing$, porque $A\subset Y$. Esto prueba la identidad y la inclusión. Como ejemplo, tomemos $X=\R$ y $Y=A=\{1/n:n\ge1\}$. Cada punto $1/n$ es aislado en $Y$; por ello $A'_Y=\varnothing$. En $\R$, el punto $0$ es de acumulación, de modo que la inclusión es estricta.
\end{solucion}

\begin{ejercicio}[Lectura 4, ejercicio 5: corrección del enunciado]
La lectura propone que la imagen de una subsucesión convergente de una sucesión acotada es cerrada en la imagen de la sucesión original. Examinar la afirmación y establecer una versión válida.
\end{ejercicio}
\begin{solucion}
La afirmación escrita es falsa. En $\R$, definamos $x_1=0$ y $x_{n+1}=1/n$. La sucesión es acotada, su imagen es $Y=\{0\}\cup\{1/n:n\ge1\}$ y la subsucesión $(x_{n+1})$ converge a $0$. Su imagen $S=\{1/n:n\ge1\}$ no es cerrada en $Y$, pues $0\in\cl S^{\,Y}\setminus S$.

La versión correcta es la siguiente: si una subsucesión converge a $\ell\in X$ y $S$ es su imagen, entonces
\[
\cl S^{\,X}=S\cup\{\ell\},\qquad
\cl S^{\,Y}=Y\cap(S\cup\{\ell\}).
\]
Para demostrar la primera igualdad, $\ell\in\cl S^{\,X}$ por convergencia y $S\subset\cl S^{\,X}$. Si $z\notin S\cup\{\ell\}$, sea $r=d(z,\ell)/3>0$. Todos los términos suficientemente avanzados están en $B(\ell,r)$, que no intersecta a $B(z,r)$. Los términos anteriores forman un conjunto finito que no contiene a $z$. Podemos reducir $r$ a un radio positivo menor que las distancias de $z$ a esos términos. La bola reducida no intersecta a $S$; así $z\notin\cl S^{\,X}$. La segunda igualdad es la fórmula de clausura relativa.

Por tanto $S$ es cerrado en $Y$ exactamente cuando $\ell\notin Y\setminus S$. En particular, lo es si $\ell\in S$ o si $\ell\notin Y$; y $S\cup\{\ell\}$ siempre es cerrado en $X$. La acotación de la sucesión original no interviene en esta versión correcta.
\end{solucion}

\begin{ejercicio}[Lectura 4, ejercicio 6]
¿Es cierto que $A\subset Y$ es abierto en $Y$ si y solo si $A=Y\cap\operatorname{int}_X A$?
\end{ejercicio}
\begin{solucion}
No. Tomemos $X=\R$, $Y=[0,1]$ y $A=Y$. El conjunto $A$ es abierto en $Y$, pero $Y\cap\operatorname{int}_{\R}A=(0,1)\ne A$.

La igualdad propuesta sí es suficiente: si $A=Y\cap\operatorname{int}_X A$, es intersección de $Y$ con un abierto de $X$. De hecho, como $\operatorname{int}_X A\subset A\subset Y$, esa igualdad equivale a $A=\operatorname{int}_X A$, es decir, a apertura en el espacio ambiente, condición más fuerte que la apertura relativa.

La fórmula general correcta para el interior relativo es
\[
\operatorname{int}_Y A
=Y\cap\operatorname{int}_X\bigl(A\cup(X\setminus Y)\bigr).
\]
Para $y\in Y$, una bola relativa está contenida en $A$ si y solo si la bola ambiental correspondiente está contenida en $A\cup(X\setminus Y)$. La equivalencia se verifica punto a punto: los puntos de la bola que no están en $Y$ no imponen restricciones relativas; los que sí están deben pertenecer a $A$.
\end{solucion}

\section{El papel del ambiente}
La clausura relativa se obtiene tomando la clausura ambiental y eliminando los puntos exteriores al subespacio. Para el interior no basta esa misma receta: primero hay que permitir todos los puntos exteriores a $Y$. Registrar el ambiente en la notación evita confundir afirmaciones sobre límites, apertura y cerradura que pueden cambiar al pasar de un espacio a otro.
